\documentclass[11pt]{article}
\usepackage[margin=1in]{geometry}
\usepackage{amsmath,amssymb,amsthm,booktabs,microtype}
\usepackage{xcolor}
\usepackage[colorlinks=true,linkcolor=blue,urlcolor=blue]{hyperref}
\allowdisplaybreaks

\newtheorem{theorem}{Theorem}
\newtheorem{proposition}[theorem]{Proposition}
\newcommand{\F}{\mathbb F}
\newcommand{\C}{\mathbb C}
\newcommand{\E}{\mathbb E}
\newcommand{\Prob}{\mathbb P}
\newcommand{\Fix}{\operatorname{Fix}}
\newcommand{\Tr}{\operatorname{Tr}}
\newcommand{\Sym}{\operatorname{Sym}}
\newcommand{\Exp}{\operatorname{Exp}}
\newcommand{\stirling}[2]{\genfrac{[}{]}{0pt}{}{#1}{#2}}

\title{The Hyperoctahedral Group on the Binary Cube\\
Fixed Points, the Full $\sigma$-MGF, and Exact Verification}
\author{Executable companion note}
\date{17 July 2026}

\begin{document}
\maketitle

\begin{abstract}
Let $B_n=C_2^n\rtimes S_n$ act affinely on $X_n=\F_2^n$ and let
$V_n=\C[X_n]$.  We prove the proposed fixed-point distribution and moment
formulas, derive the exact signed-cycle-type expression for the full
$\sigma$-MGF, and describe an independent exhaustive verification.  The
program constructs every $2^n\times2^n$ permutation matrix, computes the
target directly, traverses the relevant monomial-basis orbits, and compares
both results with the signed-cycle formula for $1\leq n\leq5$.  All checks
pass; the largest case has $|B_5|=3840$ matrices of size $32\times32$.
\end{abstract}

\section{The action and the target}
Write elements of $B_n$ as $g=(b,\pi)$, where $b\in\F_2^n$ and
$\pi\in S_n$, and set
\[
g\cdot x=b+\pi x,\qquad x\in X_n.
\]
On the basis $\{e_x:x\in X_n\}$ of $V_n$, the corresponding matrix is
\begin{equation}\label{eq:matrix}
P_g e_x=e_{g\cdot x}.
\end{equation}
Thus $P_g$ is a $2^n\times2^n$ permutation matrix.  If $c_d(g)$ denotes the
number of $d$-cycles of $g$ on the cube vertices, the generalized permutation
Molien identity gives
\begin{equation}\label{eq:molien}
\mathcal H_n(\mathbf t)
=\E_{g\in B_n}\prod_{d\geq1}Q_d(\mathbf t)^{c_d(g)},
\qquad
Q_d(\mathbf t)=\prod_{a\geq1}(1-t_a^d)^{-1}.
\end{equation}
Indeed, a $d$-cycle permutation block has determinant
$\det(I-tP_d)=1-t^d$.  Equivalently,
\[
\mathcal H_n(\mathbf t)
=\E_g\prod_{a\geq1}\det(I-t_aP_g)^{-1},
\]
which is the intended $\Exp_\sigma$ expression.

For a partition $\tau=(\tau_1,\ldots,\tau_s)$, its monomial coefficient has
the useful orbit interpretation
\begin{equation}\label{eq:orbit-coeff}
[m_\tau]\mathcal H_n
=\dim\left(\bigotimes_{j=1}^s\Sym^{\tau_j}(V_n)\right)^{B_n}
=\#\left(B_n\backslash\prod_{j=1}^s\Sym^{\tau_j}(X_n)\right).
\end{equation}
Here $\Sym^r(X_n)$ means the set of $r$-multisets of vertices.  This supplies
a direct calculation independent of determinant algebra.

\section{The fixed-point law}
Define
\[
F_n(g)=\Tr(P_g)=\#\Fix_{X_n}(g).
\]

\begin{theorem}[Exact conditional law]\label{thm:fixed-law}
For uniform $b\in\F_2^n$ and fixed $\pi\in S_n$, let $C(\pi)$ be the number
of coordinate cycles of $\pi$.  Then
\[
F_n(b,\pi)=
\begin{cases}
2^{C(\pi)},&\text{with probability }2^{-C(\pi)},\\
0,&\text{with probability }1-2^{-C(\pi)}.
\end{cases}
\]
\end{theorem}

\begin{proof}
On a coordinate cycle $(i_1,\ldots,i_\ell)$, the fixed equation
$x=b+\pi x$ becomes a cyclic system
\[
x_{i_1}=b_{i_1}+x_{i_\ell},\quad
x_{i_2}=b_{i_2}+x_{i_1},\quad\ldots\quad
x_{i_\ell}=b_{i_\ell}+x_{i_{\ell-1}}.
\]
Adding in $\F_2$ shows that it is solvable only if
$\sum_{j=1}^{\ell}b_{i_j}=0$.  This condition is also sufficient: after one
coordinate is chosen, the remaining coordinates are forced.  Consequently
each coordinate cycle imposes one independent parity condition on $b$ and,
when solvable, contributes two solutions.  There are $C(\pi)$ independent
conditions and $2^{C(\pi)}$ solutions.
\end{proof}

Let $\stirling{n}{k}$ be the unsigned Stirling number of the first kind.  Since
there are $\stirling{n}{k}$ permutations with $k$ cycles, Theorem
\ref{thm:fixed-law} yields
\begin{equation}\label{eq:distribution}
\Prob(F_n=2^k)=\frac{\stirling{n}{k}}{n!\,2^k},
\qquad 1\leq k\leq n.
\end{equation}
The remaining probability is at zero.  Using
$\sum_k\stirling{n}{k}x^k=x(x+1)\cdots(x+n-1)$ gives
\begin{equation}\label{eq:positive-probability}
\Prob(F_n>0)
=\frac{(1/2)^{\overline n}}{n!}
=\frac{\Gamma(n+\tfrac12)}{\sqrt\pi\,\Gamma(n+1)}
=\frac{\binom{2n}{n}}{4^n}
\sim\frac1{\sqrt{\pi n}}.
\end{equation}
Hence $F_n\to0$ in probability.

For every integer $m\geq1$, conditioning on $\pi$ instead gives
\[
\E_b[F_n^m\mid\pi]
=2^{-C(\pi)}2^{mC(\pi)}=2^{(m-1)C(\pi)}.
\]
Another application of the Stirling generating polynomial proves
\begin{equation}\label{eq:moments}
\boxed{\quad
\E[F_n^m]
=\frac{(2^{m-1})^{\overline n}}{n!}
=\binom{n+2^{m-1}-1}{n}.
\quad}
\end{equation}
In particular $\E F_n=1$, $\E F_n^2=n+1$,
$\E F_n^3=\binom{n+3}{3}$, and $\operatorname{Var}(F_n)=n$.  Thus convergence
in probability to zero coexists with divergent moments of every order
$m\geq2$.

There is also a direct Burnside proof of \eqref{eq:moments}.  Represent an
ordered $m$-tuple of cube vertices by an $m\times n$ binary matrix.  Coordinate
permutations permute its columns, while a flip in one cube coordinate adds
$\mathbf1=(1,\ldots,1)$ to the corresponding column.  Each column is therefore
remembered only as an element of
$\F_2^m/\langle\mathbf1\rangle$, a set of size $2^{m-1}$.  The orbit is a
multiset of $n$ such column types, so
\[
\#(B_n\backslash X_n^m)=\binom{n+2^{m-1}-1}{n}=\E[F_n^m].
\]
Combining this with \eqref{eq:orbit-coeff} gives
\begin{equation}\label{eq:moment-coeff}
[m_{(1^m)}]\mathcal H_n=\binom{n+2^{m-1}-1}{n}.
\end{equation}

\section{Signed cycle types and the full cycle index}
A signed coordinate-cycle type is a pair of sequences
$(a_\ell,b_\ell)_{\ell\geq1}$, where $a_\ell$ counts length-$\ell$
coordinate cycles with total flip parity zero and $b_\ell$ counts those with
total flip parity one.  They satisfy
\[
\sum_{\ell\geq1}\ell(a_\ell+b_\ell)=n.
\]
For a uniform element of $B_n$, the probability of this type is
\begin{equation}\label{eq:type-probability}
w(a,b)=\prod_{\ell\geq1}
\frac1{(2\ell)^{a_\ell+b_\ell}a_\ell!b_\ell!}.
\end{equation}
This follows from the usual $S_n$ cycle-type probability and the fact that
the accumulated parity on every coordinate cycle is an independent fair bit.

Let
\[
R_r(a,b)=\#\Fix_{X_n}(g^r)
\]
for any element of signed type $(a,b)$.  A positive length-$\ell$ cycle
contributes $2^{\gcd(\ell,r)}$ fixed choices.  For a negative cycle, write the
affine map as $T(x)=Rx+e_0$, with $R$ a cyclic rotation.  If
$d=\gcd(\ell,r)$, the permutation $R^r$ has $d$ coordinate orbits.  On each
orbit, the translation in $T^r$ has parity $r/d$.  The fixed equation is
therefore solvable exactly when $r/d$ is even; when solvable it has $2^d$
solutions.  Independence between coordinate cycles proves
\begin{equation}\label{eq:fixed-powers}
R_r(a,b)=
\begin{cases}
\displaystyle
2^{\sum_{\ell\geq1}(a_\ell+b_\ell)\gcd(\ell,r)},
&\displaystyle
\frac r{\gcd(\ell,r)}\text{ is even for every }\ell\text{ with }b_\ell>0,
\\[8pt]
0,&\text{otherwise}.
\end{cases}
\end{equation}

Because $R_r=\sum_{d\mid r}d\,c_d$, Möbius inversion recovers the exact
vertex-cycle counts:
\begin{equation}\label{eq:mobius}
c_d(a,b)=\frac1d\sum_{r\mid d}\mu(d/r)R_r(a,b).
\end{equation}
Substituting \eqref{eq:type-probability}--\eqref{eq:mobius} into
\eqref{eq:molien} proves the full finite-$n$ formula
\begin{equation}\label{eq:full-formula}
\boxed{
\mathcal H_n(\mathbf t)=
\sum_{\substack{(a_\ell),(b_\ell)\\
\sum_\ell\ell(a_\ell+b_\ell)=n}}
w(a,b)\prod_{d\geq1}Q_d(\mathbf t)^{c_d(a,b)}.
}
\end{equation}
This is exact, not an asymptotic or a truncation.

At total degree two, the orbit interpretation gives a transparent check.
The group is transitive on vertices, and pairs of vertices are classified by
their Hamming distance $0,1,\ldots,n$.  Therefore
\begin{equation}\label{eq:degree-two}
\mathcal H_n
=1+m_{(1)}+(n+1)\bigl(m_{(2)}+m_{(1,1)}\bigr)+\cdots.
\end{equation}
In particular, the coefficients do not stabilize as $n\to\infty$.

\section{Independent verification strategy}
The companion program avoids using the statement being tested as its only
computational route.

\paragraph{Literal matrix computation.}
For every one of the $2^n n!$ pairs $(b,\pi)$, the program builds the vertex
permutation and the matrix \eqref{eq:matrix}.  It checks that the matrices are
distinct permutation matrices.  For a partition $\tau$, it computes complete
symmetric characters from actual matrix traces using Newton's recurrence
\begin{equation}\label{eq:newton}
d\,h_d(P_g)=\sum_{r=1}^d\Tr(P_g^r)h_{d-r}(P_g),
\qquad h_0=1,
\end{equation}
and averages $\prod_jh_{\tau_j}(P_g)$ exactly.

\paragraph{Direct orbit computation.}
The monomial basis in \eqref{eq:orbit-coeff} is constructed literally.  A
breadth-first traversal using the $n$ coordinate flips and $n-1$ adjacent
coordinate transpositions counts its orbits.  This computation does not use
traces, determinants, signed types, or the proposed formula.

\paragraph{Signed-type computation.}
All nonnegative solutions $(a_\ell,b_\ell)$ of the size constraint are
generated without enumerating group elements.  Their rational weights
\eqref{eq:type-probability} are checked to sum to one.  Equations
\eqref{eq:fixed-powers} and \eqref{eq:mobius} then produce vertex cycles and
the coefficient predicted by \eqref{eq:full-formula}.

The tests also compare the entire convergent function, not only a finite
Taylor truncation.  At $(t_1,t_2)=(0.037,0.061)$ they evaluate
\[
\frac1{|B_n|}\sum_g
\frac1{\det(I-t_1P_g)\det(I-t_2P_g)}
\]
by dense determinants, by actual vertex cycles, and by the signed-type sum.

\section{Exact results}
The following coefficient vectors are ordered as
\[
\tau=(1),(2),(1,1),(3),(2,1),(1,1,1).
\]
Every entry agrees across literal matrices, direct orbit traversal, and the
signed-type formula.
\begin{center}
\begin{tabular}{@{}rrrrl@{}}
\toprule
$n$ & $|B_n|$ & $\dim V_n$ & $\Prob(F_n>0)$ & common coefficient vector\\
\midrule
1 & 2    & 2  & $1/2$   & $(1,2,2,2,3,4)$\\
2 & 8    & 4  & $3/8$   & $(1,3,3,4,7,10)$\\
3 & 48   & 8  & $5/16$  & $(1,4,4,7,13,20)$\\
4 & 384  & 16 & $35/128$& $(1,5,5,11,22,35)$\\
5 & 3840 & 32 & $63/256$& $(1,6,6,16,34,56)$\\
\bottomrule
\end{tabular}
\end{center}

For $n=1,\ldots,5$, the program additionally verifies, element by element,
the conditional parity law and reconstructs every vertex-cycle count from
\eqref{eq:fixed-powers}--\eqref{eq:mobius}.  It checks moments $m=1,2,3,4$,
giving 20 exact moment comparisons, and the six displayed coefficients,
giving 30 three-way coefficient comparisons.  In the largest numerical case,
the matrix determinant average is
\[
1.1620117926716296,
\]
and its maximum discrepancy from the cycle and signed-type evaluations is
$1.74\times10^{-14}$, consistent with floating-point roundoff.  All integer
and rational comparisons are exact.

The range $n\leq5$ is a practical exhaustive choice: it includes all 3840
elements at $n=5$ and orbit spaces as large as $|X_5|^3=32768$, while the
group order grows as $2^n n!$.  The proof and signed-type algorithm remain
valid for every $n$; only brute-force matrix enumeration is deliberately
bounded.

\section{Reproduction}
From the repository root, run
\begin{verbatim}
marimo run marimo_notebooks/hyperoctahedral_cube.py
\end{verbatim}
The notebook allows inspection of individual matrices and signed types and
shows every comparison table.  The noninteractive regression test is
\begin{verbatim}
pytest -q tests/test_hyperoctahedral_cube.py
\end{verbatim}

\section{Conclusion}
The proposed formulas are correct.  A random cube symmetry has no fixed
vertex with probability tending to one, but its rare nonzero fixed sets are
large enough that $\E F_n=1$ and all higher moments diverge.  The signed-cycle
power formula determines every vertex cycle and hence the complete
$\sigma$-MGF through \eqref{eq:full-formula}.  Exhaustive matrix, orbit, and
signed-type computations agree throughout the representative range
$1\leq n\leq5$.

\end{document}
